\documentclass[conference,a4paper,10pt]{IEEEtran}
\IEEEoverridecommandlockouts

\usepackage[T1]{fontenc}
\usepackage[utf8]{inputenc}
\usepackage{lmodern}
\usepackage{microtype}
\usepackage{textcomp}
\usepackage{xcolor}
\usepackage{booktabs}
\usepackage{array}
\usepackage{tabularx}
\usepackage{listings}
\usepackage{url}
\usepackage[hyphens]{xurl}
\usepackage{hyperref}
\hypersetup{
  colorlinks=true,
  linkcolor=black,
  citecolor=black,
  urlcolor=black,
  pdftitle={A Reproducible, Audit-Logged Container Pattern for Open-Source Intelligence Collection},
  pdfauthor={Ali Murtaza Bhutto}
}

\lstset{
  basicstyle=\ttfamily\footnotesize,
  breaklines=true,
  columns=fullflexible,
  keepspaces=true,
  frame=single,
  framerule=0.2pt,
  framesep=2pt,
  xleftmargin=4pt,
}

\begin{document}

\title{A Reproducible, Audit-Logged Container Pattern for\\Open-Source
       Intelligence Collection}

\author{\IEEEauthorblockN{Ali Murtaza Bhutto}
\IEEEauthorblockA{ORCID: 0009-0007-2787-943X\\
\texttt{alibhutto101112@gmail.com}}}

\maketitle

\begin{abstract}
Open-source intelligence (OSINT) collection in production settings is
typically performed by a fragmented toolchain that the operator
assembles by hand on a workstation. The result is poor
reproducibility, ad-hoc audit trails, and a non-trivial chance that
artefacts from one engagement leak into the next. This paper
describes \emph{docker-osint-stack}, a Compose-defined deployment
that bundles two widely used OSINT tools (SpiderFoot and theHarvester)
with a PostgreSQL audit-log and an Nginx reverse proxy into a
single-command launch, behind opinionated DevSecOps gates and an
encrypted-export script. We report measured DevSecOps numbers
(hadolint, Trivy, Dockle), measured runtime numbers (boot time, idle
RAM, smoke-test outcomes), and discuss the contribution as a
\emph{packaging} pattern rather than a new detection technique. The
artefact is open-source under MIT.
\end{abstract}

\begin{IEEEkeywords}
OSINT, Docker, container security, audit logging, reproducible
deployment, DevSecOps, SpiderFoot, theHarvester
\end{IEEEkeywords}

\section{Introduction}

OSINT collection has a tooling problem that is not, primarily, a
detection problem. The community has converged on a small number of
mature tools---SpiderFoot for aggregation across hundreds of passive
sources, theHarvester for subdomain and email enumeration---and a
sizeable body of practitioner literature on how to apply them
\cite{bazzell2023}. What is generally \emph{not} solved is the
operational shell around them: an analyst working a real engagement
needs (i) a reproducible build of every tool at a pinned version, (ii)
a tamper-evident record of who ran what against whom and under what
authorisation, (iii) a secure way to wipe and export the artefacts at
engagement close, and (iv) a defensible security posture for the
infrastructure that hosts the collection.

These needs are well understood in adjacent fields. The intelligence
literature has long argued that the auditability of an analytic chain
matters as much as the chain's content \cite{heuer1999}. Container
security guidance from NIST and CIS describes a clear set of practice
gates---non-root users, no privilege escalation, healthchecks,
minimal base images---that an OSINT deployment should meet just as
any other production container would \cite{nist800190,cisdocker}.
Yet the tooling that most practitioners use does not ship in a form
that satisfies these gates by default.

\emph{docker-osint-stack} addresses the gap by composing existing
tools, rather than building new ones, into an environment that is
trivial to stand up, hard to misuse, and \emph{measured}---in the
sense that every claim in its documentation is backed by a tool
output that the reader can regenerate.

\paragraph{Contribution.}
The contribution is a packaging pattern and the measurements that
substantiate it, not a new detection capability. Specifically:
(1) a Compose definition for a four-service stack (SpiderFoot,
theHarvester, PostgreSQL audit-log, Nginx) with the externally
reachable surface restricted to a single bound-to-localhost port;
(2) an audit-log schema (Section~\ref{sec:audit}) that records
sessions, tool invocations, normalised findings, and raw outputs in
a way that survives the destruction of any single container's
volumes; (3) a measured DevSecOps profile (Section~\ref{sec:measured})
that is honest about its non-zero CVE counts and explains the
remediation posture; and (4) a workflow-template structure that
forces an authorisation gate before any tool runs.

\section{Architecture}
\label{sec:arch}

The stack consists of four services on a private Docker bridge
network. Three are long-running (SpiderFoot, PostgreSQL, Nginx) and
one is a one-shot CLI (theHarvester) gated behind a Compose profile so
it does not start on \verb|docker compose up|.

\paragraph{External surface.}
The host port (default \texttt{127.0.0.1:8080}) is the only externally
reachable network surface. SpiderFoot's HTTP listener is bound only
on the container network, with Nginx as the lone reverse proxy. This
makes the operator-facing endpoint a single well-known location and
prevents the common failure mode where an analyst accidentally
exposes SpiderFoot's UI on a public IP.

\paragraph{Container hardening.}
Every service inherits a YAML anchor that sets
\verb|no-new-privileges:true| and a restart policy of
\verb|unless-stopped|. Both custom Dockerfiles (SpiderFoot,
theHarvester) use a two-stage build: a builder stage installs system
packages and Python wheels, and a runtime stage copies only the
required artefacts into a slim base image and runs as a dedicated
non-root user (\texttt{spider} and \texttt{harvester} respectively).
Each service that exposes a network port declares an HTTP or
\texttt{pg\_isready} healthcheck; \texttt{scripts/start.sh} polls these
healthchecks with a three-minute deadline.

\paragraph{Persistence and teardown.}
Two named volumes (\texttt{postgres-data}, \texttt{spiderfoot-data})
carry the state that should survive a container restart. The
\texttt{teardown.sh} script removes both by default; the
\texttt{--shred} flag additionally calls GNU \texttt{shred} over the
GPG-encrypted exports directory. The default-destructive behaviour is
a deliberate choice: an engagement should end in known-clean state,
and \texttt{--keep-volumes} exists for the cases where it should not.

\section{Audit-log schema}
\label{sec:audit}

The audit-log is a dedicated PostgreSQL 16 service so that destroying
a SpiderFoot or theHarvester container does not destroy the
engagement record. Schema initialisation is idempotent and loaded
from \texttt{services/audit-log/init.sql} via the standard
\texttt{/docker-entrypoint-initdb.d} mechanism in the official
PostgreSQL image \cite{postgresdocs}.

Four tables capture the audit trail:

\begin{itemize}
  \item \textbf{\texttt{investigation\_sessions}} --- one row per
    engagement or question, carrying an \texttt{operator} handle and,
    crucially, an \texttt{authorisation} reference. A \texttt{CHECK}
    constraint refuses to record an empty authorisation, encoding
    the policy that the audit-log is not a free-form notepad.
  \item \textbf{\texttt{tool\_invocations}} --- one row per
    SpiderFoot scan or theHarvester run, with the tool's arguments
    captured as JSONB so the row schema does not need to enumerate
    every possible flag.
  \item \textbf{\texttt{findings}} --- normalised entities (subdomain,
    email, IP, hash, etc.) keyed to a \texttt{finding\_severity}
    ENUM bucket (\texttt{info} through \texttt{critical}).
  \item \textbf{\texttt{raw\_outputs}} --- verbatim tool output as
    JSONB, foreign-keyed to the invocation. Storing the raw output
    alongside the normalised view lets later reanalysis run against
    the original artefact rather than against a destructive
    summary.
\end{itemize}

This structure mirrors the analytic argument
\cite{heuer1999} that intelligence work is only trustworthy when its
provenance chain---raw, processed, analysed, disseminated---is itself
inspectable. The schema makes that chain a query, not a footnote.

\section{Workflows and authorisation gate}

Four investigation templates ship with the stack: \emph{domain
reconnaissance}, \emph{social media OSINT} (organisation-centric),
\emph{credential exposure}, and a blank template. Each follows the
same six-phase structure: scope and legal gate, Priority Intelligence
Requirements (PIRs), theHarvester pass, SpiderFoot pass, triage
table, and session close. The legal-gate phase is mandatory: the
template begins with an explicit row for the authorisation reference,
and the audit-log schema's \texttt{CHECK} constraint enforces it at
storage time as well.

The intent is to make it harder for a tool run to outrun its
paperwork. This is a recurring failure mode in practitioner OSINT
\cite{bazzell2023}, and the most cost-effective place to address it
is at the template, not at the tool.

\section{Measured posture}
\label{sec:measured}

All numbers in this section are reproducible by re-running the
scripts in \texttt{tests/} and \texttt{results/}. They are not
estimates.

\subsection{Dockerfile lint (hadolint 2.12.0)}

Both custom Dockerfiles pass hadolint with zero findings under the
default ruleset.

\subsection{CVE scan (Trivy 0.70.0)}

Table~\ref{tab:trivy} reports severity counts per image. The numbers
are non-zero because SpiderFoot v4.0 (April 2022) pins dependency
versions that have accreted CVEs since release. We do not suppress
them; we report the patch path as \emph{rebuild against the same tag}
(which picks up upstream's own dependency updates when they ship) and
\emph{rebase to a newer SpiderFoot tag} once one exists. The
remediation posture treats CVE suppression as a worse failure mode
than honest reporting, in keeping with both the NIST
\cite{nist800190} and CIS \cite{cisdocker} guidance for production
container deployments.

\begin{table}[!t]
\caption{Trivy CVE counts per image (measured).}
\label{tab:trivy}
\centering
\footnotesize
\begin{tabularx}{\columnwidth}{Xrrrr}
\toprule
\textbf{Image} & \textbf{CRIT} & \textbf{HIGH} & \textbf{MED} & \textbf{LOW} \\
\midrule
\texttt{spiderfoot:v4.0}    &  8 &  46 & 120 & 118 \\
\texttt{harvester:4.10.1}   &  3 &  17 &  72 &  94 \\
\texttt{postgres:16-alpine} &  1 &  13 &  19 &   1 \\
\texttt{nginx:1.27-alpine}  &  6 &  29 &  46 &   9 \\
\bottomrule
\end{tabularx}
\end{table}

\subsection{Container best-practice (Dockle 0.4.15)}

After two documented false-positive suppressions
(see Table~\ref{tab:dockle} caption), the custom images report zero
FATAL findings; the official base images report one informational
WARN each (``last user should not be root''), which is mitigated at
Compose level by \verb|no-new-privileges| and by both nginx and
postgres themselves dropping to non-root worker processes after
init.

\begin{table}[!t]
\caption{Dockle findings per image after documented false-positive
filters. The suppressions are: \texttt{-af settings.py} for
\texttt{python-docx} (SpiderFoot) and \texttt{python-shodan}
(theHarvester), both library files literally named
\texttt{settings.py}; and \texttt{-ak KEY\_SHA512} for the
\texttt{nginx:1.27-alpine} signing-key checksum.}
\label{tab:dockle}
\centering
\footnotesize
\begin{tabularx}{\columnwidth}{Xrr}
\toprule
\textbf{Image} & \textbf{FATAL} & \textbf{WARN} \\
\midrule
\texttt{spiderfoot:v4.0}    & 0 & 0 \\
\texttt{harvester:4.10.1}   & 0 & 0 \\
\texttt{postgres:16-alpine} & 0 & 1 \\
\texttt{nginx:1.27-alpine}  & 0 & 1 \\
\bottomrule
\end{tabularx}
\end{table}

\subsection{Live boot}

A cold \verb|docker compose up| of the three long-running services
reaches the ``healthy'' state, end-to-end, in
$7.48$~seconds on a commodity laptop (Linux, 8 cores, $\approx 10$~GiB
free RAM at the time of measurement). The wall-clock figure comes
from \verb|/usr/bin/time -f '%e'| wrapping
\verb|scripts/start.sh|. Idle memory usage of the three services
totals approximately $153$~MiB
(\texttt{dos-postgres} $36.9$~MiB, \texttt{dos-spiderfoot}
$106.2$~MiB, \texttt{dos-nginx} $10.2$~MiB).

\subsection{Smoke test}

The end-to-end smoke test (\texttt{tests/test\_stack.sh}) runs ten
assertions: Compose syntax, environment file presence, full stack
boot, PostgreSQL readiness, SpiderFoot HTTP, Nginx \texttt{/healthz},
the Nginx-to-SpiderFoot proxy path, presence of the four audit-log
tables, a CTE round-trip insert across
\texttt{investigation\_sessions \(\rightarrow\) tool\_invocations
\(\rightarrow\) findings}, and a theHarvester help invocation via
the one-shot profile. All ten pass on the reference host.

\section{Threat model and limitations}

We are explicit about what the stack does \emph{not} defend against:

\begin{itemize}
  \item \textbf{TLS at the edge.} The default deployment terminates
    plain HTTP on \texttt{127.0.0.1}. Externalising the UI requires a
    TLS-terminating reverse proxy and authentication in front. We do
    not bundle either, because the choice (Caddy vs.\ Traefik vs.\
    nginx with cert-manager, basic-auth vs.\ OIDC) is engagement-
    specific.
  \item \textbf{Compromised upstream.} We pin to tags, not commit
    hashes, so a rebuild trusts whatever upstream has retroactively
    pushed behind a tag. The mitigation is to fix a commit hash in
    regulated environments. The default pins are easier to keep
    current; the more conservative pins are more reproducible.
  \item \textbf{Per-user authentication.} SpiderFoot's UI is
    single-tenant. The stack assumes one trusted operator per host.
\end{itemize}

These limitations are visible in the public README; they are not
hidden behind a marketing claim of generality.

\section{Conclusion}

\emph{docker-osint-stack} is deliberately small. It does not propose
a new OSINT methodology; the methodology is in the workflow templates
and rests on existing practitioner work \cite{bazzell2023,heuer1999}.
It does not propose a new container-security regime; the regime is
the existing NIST \cite{nist800190} and CIS \cite{cisdocker} guidance,
applied honestly. What it contributes is a packaging that makes the
combination \emph{operable in a single command} and a measurement
discipline that makes every claim in its documentation auditable.
The artefact is MIT-licensed and intended for authorised security
engagements only; the \texttt{ETHICAL\_USE.md} file enumerates the
out-of-scope cases.

\begin{thebibliography}{9}

\bibitem{bazzell2023}
M.~Bazzell, \emph{Open Source Intelligence Techniques: Resources for
Searching and Analyzing Online Information}, 10th--11th eds.,
IntelTechniques.com, 2023--2024.

\bibitem{heuer1999}
R.~J.~Heuer, \emph{Psychology of Intelligence Analysis}, Center for
the Study of Intelligence, Central Intelligence Agency, 1999.

\bibitem{nist800190}
M.~Souppaya, J.~Morello, and K.~Scarfone, ``Application Container
Security Guide,'' National Institute of Standards and Technology,
NIST Special Publication 800-190, Sep.\ 2017.

\bibitem{cisdocker}
Center for Internet Security, ``CIS Docker Benchmark,'' CIS,
Tech.\ Rep., latest version, accessed 2026.
\url{https://www.cisecurity.org/benchmark/docker}

\bibitem{postgresdocs}
PostgreSQL Global Development Group, ``PostgreSQL Documentation,''
v16, 2024--2025. \url{https://www.postgresql.org/docs/16/index.html}

\end{thebibliography}

\end{document}
